% TOP_PROBLEM: {"id": 2970, "title": "Kirby Problem 4.94", "category_id": 7, "category": {"id": 7, "name": "topology", "display_name": "Topology", "description": "Properties preserved under continuous deformations.", "slug": "topology", "order_index": 7, "created_at": "2026-07-31T15:26:25.671Z"}, "status": "open"}
% FIRST_POSTED: null
% ENTRY_KIND: statement_only
% Imported from: batch13.tex; UnsolvedMath ID 2970
\documentclass[11pt]{article}
\usepackage[margin=25mm]{geometry}
\usepackage{fontspec}
\setmainfont{Latin Modern Roman}
\usepackage{amsmath,amssymb,mathtools}
\usepackage{unicode-math}
\setmathfont{Latin Modern Math}
\usepackage{xurl,hyperref}
\hypersetup{colorlinks=true,urlcolor=blue}
\setcounter{secnumdepth}{0}
\setlength{\parindent}{0pt}
\setlength{\parskip}{5pt}
\setlength{\emergencystretch}{3em}
\newcommand{\sref}[2]{\hyperlink{source:#1:#2}{[S#2]}}
\newcommand{\cataloguescope}[1]{\paragraph{Recorded target.}#1}
\begin{document}
% UnsolvedMath ID 2970; source catalogue code KP-4.94
\setcounter{section}{2969}
\section{Kirby Problem 4.94}\label{2970}
{\small\sffamily UnsolvedMath ID 2970 · Topology}
\subsection{Definitions and mathematical statement}

\paragraph{Shared notation.}$\mathbb N=\{1,2,\ldots\}$, and $\mathbb Z,\mathbb Q,\mathbb R,\mathbb C$ denote the integers, rationals, reals and complex numbers. Any other conventions are specified locally.


A Horikawa surface is a smooth minimal complex projective surface $S$ of general type on the Noether line:
\[
c_1(S)^2=2p_g(S)-4,\qquad p_g(S)\ge3,
\]
equivalently $5c_1(S)^2=c_2(S)-36$. Here $p_g(S)=\dim_{\mathbb C}H^0(S,K_S)$, $K_S$ is the canonical line bundle, and Chern numbers are evaluations of the indicated Chern classes on the complex-oriented fundamental class. General type means Kodaira dimension two, and minimal means no holomorphic sphere of self-intersection $-1$. Deformation equivalence means being connected by a chain of smooth proper holomorphic families over connected complex bases, with the given surfaces occurring as fibers.

Equip each such $S$ with its canonical symplectic structure $\omega_S$, unique up to symplectomorphism. Its cohomology class is $c_1(K_S)$; it is obtained by normalized pluricanonical projective forms and smoothing the canonical model, as in Catanese's construction. This normalization fixes the form class and is not a choice of an arbitrary Kähler form. \sref{2970}{3}
\paragraph{Statement recorded in the supplied source.}
For every pair of homotopy-equivalent Horikawa surfaces $S,S'$ lying in different deformation classes, are their underlying smooth four-manifolds diffeomorphic? Are their canonical symplectic manifolds symplectomorphic, that is, does there exist a diffeomorphism $F:S\to S'$ with
\[
F^*\omega_{S'}=\omega_S?
\]
Both questions range over the whole indicated class of pairs, rather than only a particular odd-parameter example. \sref{2970}{2}
\subsection{Short English statement}
Are homotopy-equivalent Horikawa surfaces in different complex deformation classes nevertheless diffeomorphic, and are their canonical symplectic structures symplectomorphic?
\subsection{Sources}
\begin{description}
\item[\hypertarget{source:2970:1}{S1}] ULAM.AI and the UnsolvedMath Contributors, UnsolvedMath: A Curated Collection of Open Mathematics Problems, record 2970 (KP-4.94). CC BY 4.0. \url{https://huggingface.co/datasets/ulamai/UnsolvedMath}.
\item[\hypertarget{source:2970:2}{S2}] R. İnanç Baykur, Robion C. Kirby and Daniel Ruberman, K3—A New Problem List in Low-Dimensional Topology, Mathematical Surveys and Monographs 295 (2026), author version, Problem 4.94 and Remarks, PDF page 267. \url{https://math.berkeley.edu/sites/default/files/surv-295-ruberman-watermarked-author-pdf.pdf}.
\item[\hypertarget{source:2970:3}{S3}] Fabrizio Catanese, Canonical symplectic structures and deformations of algebraic surfaces, arXiv:math/0608110v2 (2007), Theorem 1.3, page 3, and its proof. \url{https://arxiv.org/pdf/math/0608110v2}.
\end{description}
\label{end:2970}
\end{document}
